The scanned line must consist of fields (tokens) separated by delimiters of commas
or spaces (any number of either, i.e. the delimiter pattern is [, ]+ ).

The format string is not the same as that used in scanf, sscanf or fscanf in C.
It is a string of single characters, one per conversion. Allowed
conversions are -

\begin{tabular}{| l | l |} 
\hline
'd' & integer                                   \\ \hline
'u' & unsigned integer                          \\ \hline
'x' & hexadecimal integer                       \\ \hline
'o' & octal integer                             \\ \hline
'f' & floating point value                      \\ \hline
's' & string                                    \\ \hline
'.' & convert according to the variable type    \\ \hline
'*' & skip the token                            \\ \hline
' ' & ignored                                   \\ \hline
tab & ignored                                   \\ \hline
\end{tabular}

The format string must not contain any other characters.
Spaces or tabs can be included anywhere if desired. Because the scanner cannot match explicit
fields from the format string and there are no conversion paramaters such as field widths or
precisions there is no need for a \verb"'%'" escape character. The number of actual
conversions (i.e. not including the skip) must be equal to the number of output variables.

The format argument may be null or "" in which case tokens will be converted and
assigned to the variables based on the variable types.

Tokens that are to be converted as hexadecimal values must not have a leading 0x.
Similarly a leading 0 will not be interpreted as indicating an octal value. The
conversion character alone decides which radix will be used.
